\begin{textsample}{POS Dim 4 – 2012\_c5 – Score -1.00 – 2012\_c5/enade\_2012\_tur\_30ab177814.txt}  \label{ex:f4_pos_250}
CAPÍTULO 5 DISTRIBUIÇÃO DOS CONCEITOS 5.1 PANORAMA NACIONAL DA DISTRIBUIÇÃO DOS CONCEITOS A Tabela 5.1 apresenta a quantidade e distribuição de cursos de Turismo participantes do ENADE/2012, por faixa de conceito e Grande Região.
A diferença entre os cursos tabulados neste capítulo e no capítulo 2 são os cursos sem conceito, em princípio, aqueles sem alunos concluintes que participassem da prova ou que tivessem tão somente um aluno concluinte11.
Observando -se os dados da Tabela 5.1, nota -se que, dos 191 cursos participantes, 73 ( 38,2 \% ) classificaram -se com conceito 3, o valor modal.
Este foi também o conceito modal nas regiões Nordeste ( 38,8 \% ), Sudeste ( 40,8 \% ) e Sul ( 38,9 \% ).
Nas regiões Norte e Centro-Oeste o conceito modal foi 2, com respectivamente, 36,4 \% e 47,4 \%.
O conceito 2 foi o segundo mais frequente em nível nacional ( 26,2 \%, correspondendo a 50 cursos ) e o conceito 4, o terceiro ( 15,2 \%, correspondendo a 29 cursos ).
Houve, ainda, dez cursos ( 5,2 \% ) que receberam conceito 5 e outros sete cursos ( 3,7 \% ) que receberam conceito 1.
De todos 191 cursos de Turismo, 22 ficaram sem conceito ( SC ).
A região Norte participou com 11 cursos ou 5,8 \% do total nacional.
Nesta região, o conceito 2 foi atribuído a quatro cursos dos 11 participantes, o que equivale a 36,4 \% do total regional.
Nenhum curso ficou sem conceito ( SC ).
Os demais foram avaliados com conceito 1 ( dois cursos, correspondendo a 18,2 \% ), conceito 3 ( três cursos correspondendo a 27,3 \% ), conceito 4 ( um curso correspondendo a 9,1 \% ) e conceito 5 ( também um curso ). 11 Estes cursos com somente um concluinte foram considerados no capítulo 2. 5.2 CONCEITOS POR CATEGORIA ADMINISTRATIVA E POR GRANDE REGIÃO A Tabela 5.2 apresenta a distribuição dos cursos participantes do ENADE/2012 de Turismo, por Categoria Administrativa, de acordo com os conceitos por eles alcançados, segundo as Grandes Regiões brasileiras.
Dos 191 cursos participantes, 44 ( 23,0 \% ) eram ministrados em instituições públicas e 147 ( 77,0 \% ), em privadas.
De acordo com as informações da Tabela 5.2, em termos nacionais, houve equilíbrio numérico na obtenção do conceito 5 entre instituições públicas e privadas : das dez IES que receberam este conceito, cinco eram públicas e cinco privadas.
Dos 44 cursos participantes de IES Pública, o conceito 3 foi o valor modal, atribuído a 16 cursos ( 36,4 \% ).
Entre os demais cursos de instituições públicas participantes, um obteve conceito 1 ( 2,3 \% da categoria ), dez receberam conceito 2 ( 22,7 \% ) e 12 foram alocados ao conceito 4 ( 27,3 \% ).
Como já comentado, os cinco restantes ( 11,4 \% ) receberam o conceito 5.
Nesta categoria, nenhum dos cursos ficou sem conceito ( SC ).
A região Nordeste participou com 49 cursos ou 25,7 \% do total nacional.
Como já comentado, destes, 19 cursos, 38,8 \% em termos regionais, obtiveram conceito 3, o conceito modal para a região.
Os conceitos 2 e 4 foram atribuídos a, respectivamente, dez e quatro cursos ( 20,4 \% e 8,2 \% ) e os conceitos 1 e 5 a, respectivamente, dois e três cursos ( 4,1 \% e 6,1 \% ).
Nessa região, 11 cursos ficaram sem conceito ( SC ).
Dos 76 cursos participantes da região Sudeste, como já comentado, 31 ( 40,8 \% ) obtiveram conceito 3, o conceito modal.
O conceito 1 foi atribuído a um curso ( 1,3 \% ) e o conceito 2, a 16 ( 21,1 \% ).
Outros 16 cursos receberam o conceito 4 e três outros ( 3,9 \% ) receberam o conceito 5.
Nove cursos ficaram sem conceito ( SC ).
A região Sul contou com 36 cursos ou 18,8 \%.
A predominância do conceito 3, como já comentado, foi de 38,9 \%, correspondentes a 14 dos 36 cursos participantes na região Sul.
O conceito 2 foi atribuído a 11 cursos ( 30,6 \% ) e o conceito 4, a sete cursos ( 19,4 \% ).
O conceito 5 foi recebido por três cursos ( 8,3 \% ).
Apenas um dos cursos da região Sul ficou sem conceito ( SC ).
Dos 19 cursos participantes na região Centro-Oeste, nove receberam conceito 2, o conceito modal.
Nesta região apenas um curso ficou sem conceito ( SC ).
Os demais cursos foram avaliados com conceito 1 ( dois cursos, 10,5 \% ), conceito 3 ( seis cursos, 31,6 \% ) e conceito 4 ( um curso, 5,3 \% ).
Nenhum curso tampouco foi alocado no conceito 5. 5.2 CONCEITOS POR CATEGORIA ADMINISTRATIVA E POR GRANDE REGIÃO A Tabela 5.2 apresenta a distribuição dos cursos participantes do ENADE/2012 de Turismo, por Categoria Administrativa, de acordo com os conceitos por eles alcançados, segundo as Grandes Regiões brasileiras.
Dos 191 cursos participantes, 44 ( 23,0 \% ) eram ministrados em instituições públicas e 147 ( 77,0 \% ), em privadas.
De acordo com as informações da Tabela 5.2, em termos nacionais, houve equilíbrio numérico na obtenção do conceito 5 entre instituições públicas e privadas : das dez IES que receberam este conceito, cinco eram públicas e cinco privadas.
Dos 44 cursos participantes de IES Pública, o conceito 3 foi o valor modal, atribuído a 16 cursos ( 36,4 \% ).
Entre os demais cursos de instituições públicas participantes, um obteve conceito 1 ( 2,3 \% da categoria ), dez receberam conceito 2 ( 22,7 \% ) e 12 foram alocados ao conceito 4 ( 27,3 \% ).
Como já comentado, os cinco restantes ( 11,4 \% ) receberam o conceito 5.
Nesta categoria, nenhum dos cursos ficou sem conceito ( SC ).
Na rede privada, o conceito modal também foi 3, com 57 cursos dos 147 da categoria.
Nesta categoria administrativa, 22 dos cursos ficaram sem conceito ( SC ).
Entre os demais cursos participantes, seis receberam conceito 1 e 40, conceito 2.
O conceito 4 foi atribuído a 17 cursos e o conceito 5, como já comentado, a cinco cursos.
Na análise por região, observa -se que, na região Norte, as instituições privadas participaram com oito cursos ( 72,7 \% do total regional ), dos quais nenhum ficou sem conceito ( SC ).
O conceito modal para as instituições privadas na região foi 3, com três cursos, correspondendo a 37,5 \%.
Os demais receberam os conceitos 1 ( dois cursos, 25,0 \% ), 2 ( também dois cursos ) e 4 ( um curso, 12,5 \% ).
As instituições públicas participaram com três cursos ( 27,3 \% do total regional ), dos quais nenhum ficou sem conceito ( SC ), dois receberam conceito 2 ( 66,7 \% ) e um recebeu conceito 5 ( 33,3 \% ).
Na região Nordeste, a rede privada concentrou 33 dos 49 cursos participantes, equivalentes a 67,3 \% do total da região, a menor proporção entre todas as regiões, dos quais 11 ficaram sem conceito ( SC ).
Nesta Categoria Administrativa na região, predominaram os cursos com conceito 3, 11 cursos correspondendo a exatamente 1/3 dos cursos ( 33,3 \% ) oferecidos por IES privadas no Nordeste.
Os demais foram avaliados com conceitos 1 ( dois cursos correspondendo a 6,1 \% ), 2 ( sete cursos, correspondendo a 21,2 \% ) e 5 ( também dois cursos ).
Nenhum curso recebeu o conceito 4.
As instituições públicas dessa região participaram com 16 cursos ( 32,7 \% ), dos quais oito também obtiveram conceito 3, o valor modal.
Os oito demais foram avaliados com os conceitos 2 ( três cursos ), 4 ( quatro cursos ) e 5 ( um curso ).
Nenhum curso ficou sem conceito ( SC ) nesta combinação de região e categoria.
Na região Sudeste, a proporção de cursos da rede privada, 84,2 \%, foi mais elevada do que nas demais regiões brasileiras, correspondendo a 64 dos 76 cursos participantes.
Nesta categoria, na região Sudeste, o conceito modal foi 3 ( 27 cursos ).
Nesta combinação de Categoria Administrativa e Grande Região, nove cursos não receberam conceito ( SC ).
Os demais foram avaliados com conceito 1 ( um curso ), conceito 2 ( 16 cursos ), conceito 4 ( dez cursos ) e conceito 5 ( também um curso ).
Entre os 12 cursos em instituições públicas na região Sudeste, a categoria modal foi a que recebeu conceito 4 ( seis cursos ).
Os demais foram avaliados com os conceitos 3 ( quatro cursos ) e 5 ( dois cursos ).
Na região Sudeste, nesta categoria, nenhum curso ficou sem conceito ( SC ).
As instituições privadas concentraram 28 dos 36 cursos participantes da região Sul, 77,8 \% do total regional.
O conceito modal ficou entre os conceitos 2 e 3, cada um com dez cursos.
Nesta combinação de Categoria Administrativa e Grande Região, apenas um curso ficou sem receber conceito ( SC ).
Os demais foram avaliados com conceito 4 ( cinco cursos ) e conceito 5 ( dois cursos ).
As instituições públicas na região Sul participaram com oito cursos ( 22,2 \% ), dos quais nenhum ficou sem conceito ( SC ).
O conceito modal foi 3, com quatro cursos.
Além disso, o conceito 2 foi recebido por um curso, o conceito 4 por dois cursos e o conceito 5 por um outro curso.
Nenhum curso foi alocado ao conceito 1.
Na região Centro-Oeste, 14 dos 19 cursos participantes eram de instituições privadas ( 73,7 \% em termos regionais ).
Destes, seis concentraram -se no conceito 3, conceito modal.
Nesta categoria, um curso ficou sem conceito ( SC ).
Os demais receberam conceito 1 ( um curso ), 2 ( cinco cursos ) e 4 ( também um curso ).
Nenhum curso foi alocado ao conceito 5.
Dos cinco cursos de instituições públicas, quatro foram avaliados no conceito 2, conceito modal, e o outro recebeu conceito 1.
Nesta região, nenhum curso de IES pública ficou sem conceito ( SC ). 5.3 CONCEITOS POR ORGANIZAÇÃO ACADÊMICA E POR GRANDE REGIÃO Na Tabela 5.3 encontra -se a distribuição dos conceitos atribuídos aos cursos participantes do ENADE/2012 na Área de Turismo, por Organização Acadêmica, segundo as Grandes Regiões brasileiras.
Dos 191 cursos participantes, 81 eram oferecidos em Universidades, 24 em Centros Universitários e os demais 86 em Faculdades.
Esta distribuição corresponde a, respectivamente, 42,4 \%, 12,6 \% e 45,0 \% dos cursos.
De acordo com os dados apresentados, de todos os dez cursos avaliados com conceito 5, sete eram vinculados a Universidades.
Apenas três cursos em Universidades ficaram sem conceito ( SC ).
Esse tipo de Organização Acadêmica teve o conceito 4 como modal, com 26 cursos ( 32,1 \% ).
Os demais cursos de Universidades avaliados receberam os conceitos 1 ( dois cursos ), 2 ( 18 cursos ), 3 ( 25 cursos ) e 5 ( sete cursos, como já mencionado ).
Entre os cursos em Centros Universitários, o conceito modal foi 3, com 12 cursos ( 50,0 \% ).
Neste tipo de organização acadêmica cinco cursos ficaram sem conceito ( SC ) e os outros sete cursos, neste tipo de Organização Acadêmica, receberam o conceito 2.
Nas Faculdades, 14 dos 84 cursos ficou sem conceito ( SC ) e 36 ( 41,9 \% ) receberam o mesmo conceito modal dos Centros Universitários, conceito 3.
Dos demais cursos neste tipo de Organização Acadêmica, cinco receberam conceito 1, 25 conceito 2, três conceito 4 e outros três cursos, o conceito 5.
Considerando -se separadamente as regiões brasileiras, verifica -se que, na região Norte, as Universidades concentraram quatro dos 11 cursos participantes.
Dos cursos em Universidades, nenhum ficou sem conceito ( SC ) ou recebeu conceitos 1 ou 3.
Dois cursos receberam conceito 2 e os demais dois ficaram distribuídos igualmente nos conceitos 4 e 5 ( um curso cada ).
Os Centros Universitários da região Norte foram representados por dois cursos, ambos receberam conceito 3.
As Faculdades participaram com cinco cursos na região Norte.
Nenhum ficou sem conceito ( SC ) ou recebeu os conceitos 4 ou 5.
Estes receberam conceito 1 ( dois cursos ), 2 ( também dois cursos ) e conceito 3 ( um curso ).
Na região Nordeste, as Universidades participaram com 18 dos 49 cursos na Área de Turismo da região.
Apenas um dos cursos em Universidades no Nordeste ficou sem conceito ( SC ), enquanto nenhum recebeu conceito 1.
O conceito modal foi 3, com nove cursos.
Os demais receberam conceito 2 ( três cursos ), 4 ( quatro cursos ) e conceito 5 ( um curso ).
Os Centros Universitários contaram com três cursos participantes na região Nordeste, dois deles com conceito 3 e um sem conceito ( SC ).
As Faculdades foram representadas por 28 cursos na região Nordeste.
O conceito modal foi 3, com oito cursos.
Nesta combinação de organização acadêmica e região, nove cursos ficaram sem conceito ( SC ).
Dos cursos restantes, dois receberam conceito 1, sete receberam conceito 2 e outros dois receberam conceito 5.
Na região Sudeste, as Universidades concentraram 37 dos 76 cursos da região.
Entre os cursos em Universidades na região, o conceito modal foi 4 com 14 cursos, e dois cursos ficaram sem conceito ( SC ).
Os demais cursos receberam os conceitos 1 ( um curso ), 2 ( seis cursos ), 3 ( 11 cursos ) e 5 ( três cursos ).
Os Centros Universitários participaram com 12 cursos na região Sudeste, dos quais cinco obtiveram conceito modal, 2, e quatro ficaram sem conceito ( SC ).
Os demais três receberam o conceito 3.
As Faculdades foram representadas por 27 cursos na região Sudeste, que se distribuíram nos conceitos 2 ( cinco cursos ), 3 ( 17 cursos, conceito modal ) e 4 ( dois cursos ), enquanto três deles ficaram sem conceito ( SC ).
Dos 36 cursos da região Sul, 16 eram de Universidades, para os quais o conceito modal foi 4, com seis cursos.
Nesse tipo de organização, nenhum dos cursos ficou sem conceito ( SC ) e os demais receberam os conceitos 2 ( três cursos ), 3 ( cinco cursos ) e 5 ( dois cursos ).
Os Centros Universitários da região Sul tiveram cinco cursos participantes, os quais um recebeu conceito 2 e quatro receberam conceito 3.
Foram 15 os cursos vinculados a Faculdades na região Sul, e apenas um destes ficou sem conceito ( SC ).
Os demais cursos receberam os conceitos 2 ( sete cursos, conceito modal ), 3 ( cinco cursos ), 4 ( um curso ) e 5 ( também um curso ).
Na região Centro-Oeste, seis dos 19 cursos eram de Universidades.
Nesse tipo de organização, nenhum curso ficou sem conceito ( SC ) e o conceito modal foi 2, com quatro cursos.
Os outros dois cursos obtiveram os conceitos 1 e 4 ( um curso cada ).
Os Centros Universitários da região Centro-Oeste contaram com dois cursos, dos quais um recebeu o conceito 2 e o outro, o conceito 3.
Nenhum curso ficou sem conceito ( SC ).
Dos 11 cursos em Faculdades na região Centro-Oeste, apenas um ficou sem conceito ( SC ).
Dos demais, um recebeu conceito 1, quatro receberam conceito 2 e cinco receberam o conceito 3.

% matched lemmas: 
\end{textsample}
